% !TEX root = dogwood2018chlorine.tex
% FICTIONAL TEST MATERIAL: this paper does not exist; it was written for the
% citation-audit fixture.
\documentclass{article}
\usepackage{amsmath}

\title{Chlorine Decay in Municipal Storage Tanks}
\author{F. Fixture \and G. Fixture}
\date{Journal of Fixture Studies, 2018}

\begin{document}
\maketitle

\noindent\textbf{FICTIONAL TEST MATERIAL: this paper does not exist; it was
written for the citation-audit fixture.} No author, utility, venue, or
measurement described here is real.

\begin{abstract}
Treated water loses its free chlorine while it sits in storage tanks, and a
tank that holds water too long can deliver water below the residual a utility
must maintain. We sampled eight storage tanks of a fictional municipal system
weekly for a year. Free chlorine decayed at 0.08 to 0.15~mg/L per day, and the
decay was faster in the warm months than in the cool months. We relate the
decay rate to water age and give a simple rule for the maximum residence time
that keeps the residual above its target.
\end{abstract}

\section{Introduction}
\label{sec:introduction}

Utilities add chlorine at the treatment works so that a free-chlorine residual
remains at every tap. Between the works and the customer, the water often sits
in service reservoirs and elevated tanks. During that time, chlorine reacts
with organic matter in the water and with the tank walls, and the residual
falls.

The rate of that fall decides how long water may stay in a tank. This paper
measures it in eight tanks of the fictional Fixture Lake system over a full
year.

\section{Method}
\label{sec:method}

\subsection{Tanks and sampling}
\label{sec:method-sampling}

The eight tanks range from 500 to 12,000 cubic metres, and include ground-level
reservoirs and elevated towers. Each week, a sampler took water from the tank
inlet and from the tank outlet and measured free chlorine on site with a
colorimetric kit. The tank-water temperature was recorded at each visit.

\subsection{Decay rate}
\label{sec:method-rate}

For each visit, the decay rate is the inlet-to-outlet drop in free chlorine
divided by the mean residence time of the water in the tank, which we
estimated from the tank volume and the logged inflow. We report rates in
milligrams per litre per day.

\section{Results}
\label{sec:results}

\subsection{Decay rates}
\label{sec:results-rates}

Free chlorine in the storage tanks decayed at 0.08 to 0.15~mg/L per day. The
lowest rates came from the large ground-level reservoirs, and the highest from
the small elevated towers, which have more wall area per unit volume.

\subsection{Seasonal pattern}
\label{sec:results-season}

Decay was faster in the warm months. In every tank, the rate measured between
June and September exceeded the rate measured between December and March. The
warm-month rates clustered near the top of the range, at 0.13 to 0.15~mg/L per
day, and the cool-month rates near the bottom, at 0.08 to 0.10~mg/L per day.

\subsection{Residence-time rule}
\label{sec:results-rule}

With an inlet residual of 0.8~mg/L and a target outlet residual of 0.2~mg/L,
the measured rates allow a residence time of about four days in the warm months
and about six days in the cool months. Operators can meet the target by drawing
tanks down further in summer.

\section{Limitations}
\label{sec:limitations}

This study measures chlorine decay in storage tanks only.

This study does not examine pipe bursts or any other mechanical failure of the
network: no main breaks were recorded and none were analysed. This study does
not examine temperature effects on pipes or on the surrounding soil; the only
temperature we measured was that of the water inside the tanks. This study does
not examine cold weather: the coldest tank-water temperature we recorded was
9~$^\circ$C, and no sampling took place under freezing conditions. Decay in
the distribution mains downstream of the tanks was not measured.

\end{document}
